\documentclass[10pt, aspectratio=169]{beamer}
% Preámbulo modular para presentaciones Beamer (Modelación Estadística)
% Diseño y paleta institucional del Tecnológico de Monterrey

\documentclass[aspectratio=169,xcolor={dvipsnames,table}]{beamer}

\usepackage[utf8]{inputenc}
\usepackage[spanish,mexico]{babel}
\usepackage{amsmath,amssymb,amsfonts,mathtools}
\usepackage{graphicx}
\usepackage{listings}
\usepackage{booktabs}
\usepackage{tikz}
\usetikzlibrary{matrix,arrows,decorations.pathmorphing,shapes.geometric,positioning}

% Paleta Institucional Tecnológico de Monterrey
\definecolor{TecAzul}{HTML}{0039A6}       % Azul TEC: color primario institucional
\definecolor{TecAzulOscuro}{HTML}{1A2E51} % Azul marino oscuro
\definecolor{TecRojo}{HTML}{EC2661}       % Rojo de acento (paleta miTec)
\definecolor{TecGris}{HTML}{646464}       % Gris neutro para comentarios y subtítulos
\definecolor{TecFondoBloque}{HTML}{F4F6F9}% Fondo suave para bloques

% Configuración de Tema Beamer
\usetheme{Boadilla}
\usecolortheme[named=TecAzulOscuro]{structure}
\setbeamercolor{alerted text}{fg=TecRojo}
\setbeamercolor{palette primary}{bg=TecAzulOscuro,fg=white}
\setbeamercolor{palette secondary}{bg=TecAzul,fg=white}
\setbeamercolor{palette tertiary}{bg=TecAzulOscuro!80!black,fg=white}
\setbeamercolor{title}{fg=TecAzulOscuro,bg=TecFondoBloque}
\setbeamercolor{frametitle}{fg=TecAzulOscuro,bg=TecFondoBloque}

% Bloques personalizados elegantes
\setbeamercolor{block title}{bg=TecAzulOscuro,fg=white}
\setbeamercolor{block body}{bg=TecFondoBloque,fg=black}
\setbeamercolor{block title alerted}{bg=TecRojo,fg=white}
\setbeamercolor{block body alerted}{bg=TecRojo!8,fg=black}
\setbeamercolor{block title example}{bg=TecAzul,fg=white}
\setbeamercolor{block body example}{bg=TecAzul!8,fg=black}

% Redondear bordes y añadir sombra sutil
\setbeamertemplate{blocks}[rounded][shadow=true]
\setbeamertemplate{navigation symbols}{}

% Pie de página limpio con número de diapositiva
\setbeamertemplate{footline}{
  \leavevmode%
  \hbox{%
  \begin{beamercolorbox}[wd=.7\paperwidth,ht=2.25ex,dp=1ex,left]{author in head/foot}%
    \hspace*{2ex}\insertshortauthor~~(\insertshortinstitute) \hfill \insertshorttitle
  \end{beamercolorbox}%
  \begin{beamercolorbox}[wd=.3\paperwidth,ht=2.25ex,dp=1ex,right]{date in head/foot}%
    \insertshortdate{}\hspace*{2em}
    \insertframenumber{} / \inserttotalframenumber\hspace*{2ex}
  \end{beamercolorbox}}%
  \vskip0pt%
}

% Configuración de Listings de Python para visualización pedagógica de código
\lstset{
    language=Python,
    basicstyle=\ttfamily\footnotesize,
    keywordstyle=\color{TecAzulOscuro}\bfseries,
    stringstyle=\color{TecRojo},
    commentstyle=\color{TecGris}\itshape,
    showstringspaces=false,
    breaklines=true,
    frame=lines,
    rulecolor=\color{TecAzulOscuro!30},
    backgroundcolor=\color{TecFondoBloque},
    aboveskip=0.5em,
    belowskip=0.5em,
    literate={á}{{\'a}}1 {é}{{\'e}}1 {í}{{\'i}}1 {ó}{{\'o}}1 {ú}{{\'u}}1
             {Á}{{\'A}}1 {É}{{\'E}}1 {Í}{{\'I}}1 {Ó}{{\'O}}1 {Ú}{{\'U}}1
             {ñ}{{\~n}}1 {Ñ}{{\~N}}1 {¿}{{?`}}1 {¡}{{!`}}1
}

% Metadatos institucionales por defecto
\title[Modelación Estadística]{Modelación Estadística para Ciencia de Datos e Ingeniería}
\author[J. Castillo Colmenares]{Juliho David Castillo Colmenares}
\institute[Tec de Monterrey]{Tecnológico de Monterrey \\ Escuela de Ingeniería y Ciencias}
\date{\today}

% Comandos y macros estadísticas compartidas para las presentaciones Beamer (Metropolis)

% Conjuntos numéricos básicos
\newcommand{\R}{\mathbb{R}}
\newcommand{\N}{\mathbb{N}}
\newcommand{\Q}{\mathbb{Q}}
\newcommand{\Z}{\mathbb{Z}}
\newcommand{\set}[1]{\left\{ #1 \right\}}
\newcommand{\abs}[1]{\left|#1\right|}
\newcommand{\norm}[1]{\left\| #1 \right\|}

% Comandos estadísticos y probabilísticos del curso
\newcommand{\Var}{\operatorname{Var}}
\newcommand{\cov}{\operatorname{Cov}}
\newcommand{\corr}{\rho}
\newcommand{\comb}[2]{\begin{pmatrix} #1 \\ #2 \end{pmatrix}}
\newcommand{\s}{\sigma}
\newcommand{\particion}{\mathfrak{P}}
\newcommand{\card}[1]{\#\left( #1 \right)}
\newcommand{\seq}[3]{\set{#1}_{#2}^{#3}}
\newcommand{\E}{\mathbb{E}}
\newcommand{\Prob}{\mathbb{P}}

% Entornos pedagógicos y cajas en español académico natural
\newenvironment{discusion}{%
  \begin{alertblock}{Pregunta de Discusión}%
}{%
  \end{alertblock}%
}

\newenvironment{reflexion}{%
  \begin{alertblock}{Reflexión para el Aula}%
}{%
  \end{alertblock}%
}

\newenvironment{paradoja}{%
  \begin{alertblock}{Paradoja de Exploración}%
}{%
  \end{alertblock}%
}

\newenvironment{motivacion}{%
  \begin{exampleblock}{Motivación: Ciencia de Datos en la Práctica}%
}{%
  \end{exampleblock}%
}

\newenvironment{retoclase}{%
  \begin{block}{Reto Rápido en Equipo}%
}{%
  \end{block}%
}


\title[Fundamentos de Docimasia]{Sección 07.01: Fundamentos de Pruebas de Hipótesis}
\subtitle{$H_0$ vs. $H_1$, Errores Tipo I y II, y Potencia Estadística}
\author[J. Castillo Colmenares]{Juliho Castillo Colmenares}
\institute{Tecnológico de Monterrey}
\date{\vspace{-1.2cm}}

\begin{document}

% Slide 1: Portada
\begin{frame}[plain]
  \vspace{-0.8cm}
  \titlepage
\end{frame}

% Slide 2: Hoja de Ruta
\begin{frame}{Hoja de Ruta --- Capítulo 07: Docimasia (Pruebas de Hipótesis)}
  \footnotesize
  ¿Dónde nos encontramos en el desarrollo modular del capítulo? \pause
  \begin{itemize}\itemsep=0.08em
    \item \textbf{07.01} Fundamentos: $H_0$ vs. $H_1$, Errores Tipo I y II, y Potencia \emph{(Hoy)} \pause
    \item \textbf{07.02} Pruebas $Z$ y $t$ para Medias de Una y Dos Muestras \pause
    \item \textbf{07.03} Pruebas de Bondad de Ajuste $\chi^2$ \pause
    \item \textbf{07.04} Tablas de Contingencia y Pruebas de Independencia
  \end{itemize}
  \vspace{-0.05cm}
  \begin{block}{Objetivo de la Sesión}
    Formalizar el marco lógico y probabilístico de toda prueba de hipótesis: cómo se plantean $H_0$ y $H_a$, qué son los errores Tipo I ($\alpha$) y Tipo II ($\beta$), qué mide la potencia de una prueba, y cómo se determina el tamaño de muestra necesario para alcanzar una potencia objetivo.
  \end{block}
\end{frame}

% Slide 3: Motivación
\begin{frame}{De la Estimación a la Decisión}
  \footnotesize
  \begin{motivacion}
    En el Capítulo 06 aprendimos a \emph{estimar} parámetros poblacionales mediante puntos e intervalos. Ahora damos el siguiente paso lógico: usar los datos muestrales para \emph{decidir} entre dos afirmaciones contrapuestas sobre esos parámetros --- el corazón de la inferencia estadística aplicada.
  \end{motivacion}

  \vspace{0.15cm}
  \pause
  \begin{itemize}\itemsep=0.1cm
    \item \textbf{Hipótesis nula ($H_0$):} la afirmación vigente, que se asume verdadera hasta tener evidencia en contra. \pause
    \item \textbf{Hipótesis alternativa ($H_a$):} lo que buscamos evidenciar; niega a $H_0$. \pause
    \item \textbf{Toda decisión es incierta:} ninguna prueba es infalible; debemos cuantificar el riesgo de equivocarnos.
  \end{itemize}
\end{frame}

% Slide 4: Notación formal y tipos de prueba
\begin{frame}{Notación Formal y Tipos de Prueba}
  \footnotesize
  \begin{block}{Formulación de Hipótesis}
    $$H_0: \mu = \mu_0 \qquad \text{contra} \qquad H_a: \mu \neq \mu_0 \ \text{(o } < \ \text{o } > \text{)}$$
  \end{block}
  \pause

  \vspace{0.1cm}
  \begin{alertblock}{Tres Estructuras Posibles}
    \begin{itemize}\itemsep=0.05cm
      \item \textbf{Cola izquierda:} $H_a: \mu < \mu_0$.
      \item \textbf{Cola derecha:} $H_a: \mu > \mu_0$.
      \item \textbf{Dos colas:} $H_a: \mu \neq \mu_0$.
    \end{itemize}
    La dirección de $H_a$ siempre refleja la sospecha o afirmación que se busca contrastar.
  \end{alertblock}
\end{frame}

% Slide 5: Errores Tipo I y II
\begin{frame}{Los Dos Tipos de Error en la Docimasia}
  \footnotesize
  Toda decisión estadística basada en una muestra puede equivocarse de dos maneras distintas: \pause

  \vspace{0.1cm}
  \begin{block}{Error Tipo I ($\alpha$)}
    Rechazar $H_0$ cuando en realidad es verdadera.
    $$\alpha = P(\text{Rechazar } H_0 \mid H_0 \text{ verdadera})$$
  \end{block}
  \pause

  \vspace{0.1cm}
  \begin{alertblock}{Error Tipo II ($\beta$)}
    No rechazar $H_0$ cuando en realidad es falsa.
    $$\beta = P(\text{No rechazar } H_0 \mid H_0 \text{ falsa})$$
  \end{alertblock}
\end{frame}

% Slide 6: Potencia y su fórmula de tamaño de muestra
\begin{frame}{Potencia y Tamaño de Muestra}
  \footnotesize
  \begin{block}{Potencia de la Prueba}
    $$\text{Potencia} = 1-\beta = P(\text{Rechazar } H_0 \mid H_0 \text{ falsa})$$
    Es la capacidad de la prueba para detectar un efecto real cuando en verdad existe.
  \end{block}
  \pause

  \vspace{0.1cm}
  \begin{alertblock}{Tamaño de Muestra para una Potencia Objetivo (Prueba $Z$, una cola)}
    $$n = \left(\frac{(Z_\alpha+Z_\beta)\sigma}{\mu_a-\mu_0}\right)^2$$
    A mayor potencia exigida o menor $\alpha$ tolerado, mayor es el $n$ requerido; a mayor tamaño de efecto ($\mu_a-\mu_0$), menor el $n$ necesario.
  \end{alertblock}
\end{frame}

% Slide 7: Nivel de significación, valor p y regla de decisión
\begin{frame}{Significación, Valor $p$ y Regla de Decisión}
  \footnotesize
  \begin{itemize}\itemsep=0.12cm
    \item \textbf{Nivel de significación ($\alpha$):} la probabilidad máxima aceptable de Error Tipo I, fijada \emph{antes} de observar los datos (típicamente $0.10$, $0.05$ o $0.01$). \pause
    \item \textbf{Valor $p$:} la probabilidad de observar un resultado al menos tan extremo como el obtenido, asumiendo que $H_0$ es verdadera. \pause
    \item \textbf{Regla de decisión:} si $p\text{-valor} \le \alpha$, se rechaza $H_0$; en caso contrario, no se rechaza (lo cual \emph{no} equivale a aceptarla).
  \end{itemize}
\end{frame}

% Slide 8: Lab Python (Bloque 1)
\begin{frame}[fragile]{Laboratorio en Python: Tasa de Error Tipo I}
  \scriptsize
  Verificación Monte Carlo de que la tasa de falso rechazo iguala a $\alpha$ bajo $H_0$ verdadera (\texttt{07.01\_hypothesis\_testing\_basics.py}):
  {\lstset{basicstyle=\fontsize{4pt}{4.8pt}\selectfont\ttfamily,aboveskip=0.1em,belowskip=0.1em}
  \lstinputlisting[firstline=17, lastline=36]{../../code/07_pruebas_hipotesis/07.01_hypothesis_testing_basics.py}}
\end{frame}

% Slide 9: Lab Python (Bloque 2)
\begin{frame}[fragile]{Laboratorio en Python: Error Tipo II y Función de Potencia}
  \scriptsize
  Cálculo analítico y verificación por simulación de $\beta$ y la potencia:
  {\lstset{basicstyle=\fontsize{4pt}{4.8pt}\selectfont\ttfamily,aboveskip=0.1em,belowskip=0.1em}
  \lstinputlisting[firstline=39, lastline=62]{../../code/07_pruebas_hipotesis/07.01_hypothesis_testing_basics.py}}
\end{frame}

% Slide 10: Lab Python (Bloque 3)
\begin{frame}[fragile]{Laboratorio en Python: Tamaño de Muestra para Potencia Objetivo}
  \scriptsize
  Cálculo del $n$ mínimo y verificación empírica de la potencia lograda:
  {\lstset{basicstyle=\fontsize{4pt}{4.8pt}\selectfont\ttfamily,aboveskip=0.1em,belowskip=0.1em}
  \lstinputlisting[firstline=70, lastline=92]{../../code/07_pruebas_hipotesis/07.01_hypothesis_testing_basics.py}}
\end{frame}

% Slide 11: Salida Terminal
\begin{frame}[fragile]{Laboratorio en Python: Salida en Terminal}
  \scriptsize
  Salida estándar generada al ejecutar el script del laboratorio en Python:
  \vspace{0.02cm}
  \begin{block}{Salida Estándar en Consola}
    \fontsize{4pt}{5pt}\selectfont
    \begin{verbatim}
=== Block 1: Type I Error Rate ===
Critical value z_crit (alpha=0.05): 1.6449
Empirical Type I error rate: 0.0503  (nominal alpha: 0.0500)

=== Block 2: Type II Error and Power ===
Analytic beta (mu_a=108): 0.1534, Analytic power: 0.8466
Empirical power via Monte Carlo: 0.8460

=== Block 3: Sample Size for Target Power ===
Exact n = 30.11 -> rounded up: n=31
Verification with n=31: empirical power = 0.9073 (target=0.90)
    \end{verbatim}
  \end{block}
\end{frame}

% Slide 16: Cierre
\begin{frame}{Cierre de Sección: Fundamentos de Docimasia}
  \begin{reflexion}
    Toda prueba de hipótesis es un procedimiento de decisión bajo incertidumbre, calibrado explícitamente por dos riesgos contrapuestos: $\alpha$ (rechazar de más) y $\beta$ (rechazar de menos). Diseñar bien un experimento significa elegir $\alpha$, $n$ y la potencia deseada de manera consciente, no accidental.
  \end{reflexion}

  \vspace{0.2cm}
  \pause
  \begin{block}{Lecciones Clave}
    \begin{itemize}\itemsep=0.08cm
      \item \textbf{$H_0$ vs. $H_a$:} el valor conocido/afirmado va en $H_0$; la sospecha o hallazgo va en $H_a$.
      \item \textbf{$\alpha$ y $\beta$:} riesgos complementarios; reducir uno (con $n$ fijo) generalmente aumenta el otro.
      \item \textbf{Potencia y tamaño de muestra:} $n=((Z_\alpha+Z_\beta)\sigma/(\mu_a-\mu_0))^2$ vincula el diseño experimental con el riesgo tolerado.
    \end{itemize}
  \end{block}
\end{frame}

% Slide 17: Perspectiva Modular
\begin{frame}{Perspectiva Modular --- El Próximo Reto}
  \scriptsize
  \begin{retoclase}
    Ya dominamos el marco lógico completo de la docimasia. La Sección 07.02 aplicará este marco a las pruebas $Z$ y $t$ para una y dos medias --- muestras independientes con varianzas iguales o distintas (Welch), y muestras pareadas --- extendiendo el ejemplo de una sola media a comparaciones entre grupos.
  \end{retoclase}

  \vspace{0.08cm}
  \pause
  \begin{alertblock}{Preparación Recomendada}
    \begin{itemize}\itemsep=0.06cm
      \item Repasa la distribución $t$ de Student y sus grados de libertad (Sección 05.04).
      \item Reflexiona sobre por qué combinar varianzas ($S_p^2$) requiere el supuesto de homoscedasticidad.
    \end{itemize}
  \end{alertblock}
\end{frame}

\end{document}
